\section{Twist-2 operators}
QCD factorization relates the moments of the structure functions
and the moments of the parton distribution functions.
Taking the example of $F_1$, the relation is given by
e.g.
\be
\int_0^1 dx~x^{n-1} F_1(x,Q^2) = \sum_a C_{1,n}^a(Q^2/\mu^2)
A_n^{a/h}(\mu)
\label{eq:def_mom}
\ee
where $A_n^{a/h}(\mu)$ represents the reduced matrix element
of certain local operators, denoted by $O_n$, between the hadron states, denoted as $h$, and
$C_{1,n}^a(Q^2/\mu^2)$ are the corresponding Wilson coefficients, calculable
in perturbation theory.
In the context of DIS kinematics,
the leading contribution, given in Eq.~\eqref{eq:def_mom},
arises from local operators
with the lowest twist, $\tau = d_O - n = 2$, where $d_O$ is the physical
dimension of the local operator and $n$ is its spin.
In the parton model the reduced matrix element is related to the moments
of the PDFs, $A_n^{a/h} = \left\langle x^{n-1} \right \rangle_{a/h}$~\cite{Curci:1980uw,Collins:1981uw}.

The form of the local operators $O_n$ can be constrained
by requiring that they belong
to irreducible representations of the Lorentz group.
This guarantees that the renormalization procedure does not generate
mixing with operators of the same or lower dimensions.
Lattice QCD simulations are performed in a Euclidean setup though, and
the symmetry group in a Euclidean lattice is the hyper-cubic group H($4$).
The reduced symmetry does not protect the operators from a complicated mixing pattern,
that gets more complicated for larger $n$.
Once the continuum limit has been performed, one recovers O($4$) symmetry.

The construction of irreducible representations of O($4$) symmetry group
is standard material (see, for example, Ref.~\cite{hamermesh1962group}).
For the purpose of this work, the only relevant result is that
traceless and symmetrized rank-$n$ tensors are an irreducible
representation of O($4$).

We restrict ourselves to unpolarized structure functions,
meaning that we consider hadronic matrix elements averaged
over the hadron polarization.
Other moments of different distributions can be studied
using the same procedure described here.
We further limit our consideration to flavor non-singlet
operators to avoid mixing with gluonic operators,
but the same method is applicable to singlet operators.

We consider the following set of operators
\be
O_n^{rs}(x) = O^{rs}_{\mu_1 \cdots \mu_n}(x) =
\psibar^r(x) \gmuopen1 \lrDmu2 \cdots \lrDmuclose{n} \psi^s(x)\,,
\label{eq:t2}
\ee
with quark fields, $\psi^r$, of different flavors, $r \neq s$, and the
covariant derivative $D_\mu = \partial_\mu + G_\mu$, where $G_\mu$
represent the gluon field.
The symmetrization of the covariant derivative
$\lrD_\mu = \frac{1}{2}\left({\overset\rightarrow{D}_\mu} -
{\overset\leftarrow{D}_\mu}\right)$
guarantees a definite transformation under charge conjugation.

In the continuum,
traceless non-singlet twist-2 operators, with symmetrized
Lorentz indices,~\footnote{We denote the traceless operators with $\widehat{O}_n$.}
renormalize multiplicatively.
For example, in the $\MS$ scheme, the bare operator $\widehat{O}_{n,{\text{B}}}^{rs}$
renormalizes as follows
\be
\widehat{O}_n^{rs}(x) = Z^{\MS}_n \widehat{O}_{n,{\text{B}}}^{rs}(x)\,,
\ee
with renormalization constant $Z^{\MS}_n$.
The intrinsic non-perturbative nature of the hadronic matrix elements
suggests the use of lattice QCD. However, the breaking
of the rotational group symmetry into the hypercubic group H($4$)
adds complications to the renormalization.

On an hypercubic lattice, the Lorentz indices of the
twist-2 operators must be chosen while considering the reduced
hypercubic symmetry group H($4$).
Irreducible representations of O($4$) generally become reducible
representations of H($4$) inducing
unwanted mixings under
renormalization~\cite{Kronfeld:1984zv,Martinelli:1987zd,Beccarini:1995iv,Gockeler:1996mu}.
The irreducible representations of H($4$)
allow mixings with lower-dimensional operators,
and, additionally, complicate the mixing with operators of the same dimensions.
A notable example of this complication occurs with the operator $O_3$.
In the continuum, after
symmetrizing over the Lorentz indices and subtracting the traces,
the renormalization is multiplicative (e.g. in the $\MS$ scheme).
However, on the lattice, terms like $1/a^2 \delta_{\mu_i \mu_j} \cos( a p_{\mu_j})$
are allowed by H($4$) symmetry and cannot be removed even after subtracting
the trace~\cite{Kronfeld:1984zv}.
This pertains to the fact that $O_3$ belongs
to different irreducible representations of H($4$)
depending on the choice of the Lorentz indices.
The operator $O_3$ with
all the indices equal, i.e. $\mu_1 = \mu_2 = \mu_3$, mixes with the
lower-dimensional vector current, causing a power divergence proportional
to $1/a^2$. For the operators with $\mu_1 \neq \mu_2 = \mu_3$,
such as the irreducible representation $O_{411} - O_{433}$,
the renormalization is multiplicative, up to a small mixing with
a different operator of the same dimension.\footnote{
    The operator $O_{411} - O_{433}$ belongs to an irreducible
    representation of H($4$) that appears more than once.
    Therefore, a mixing between operators
    belonging to these equivalent representations can occur~\cite{Beccarini:1995iv}.
    In perturbation theory the mixing seems to be numerically small.
}
We note that while $O_{411}$ is also affected by power divergences,
the subtraction $O_{411} - O_{433}$
guarantees their removal.
Another choice is $\mu_1 \neq \mu_2 \neq \mu_3$. In this case, the operator belongs
to another irreducible representation of H($4$) and renormalizes multiplicatively
without introducing power divergences.
However, this last example, while optimal in terms of renormalization, is far from
optimal when calculating hadronic matrix elements\footnote{In Eq.~\eqref{eq:matrix} for simplicity
we omit the flavor indices.}
\be
\left \langle h(p)| \widehat{O}_n |h(p) \right\rangle\ =
p_{\mu_1} \cdots p_{\mu_n} A_n^{h}(\mu)\,,
\label{eq:matrix}
\ee
because it requires nonvanishing external momentum in at least $2$ spatial directions.
This substantially degrades the quality of the signal-to-noise ratio
that is realized in numerical lattice QCD simulations.
For $O_{411} - O_{433}$, one needs a nonvanishing external momentum in at least
$1$ spatial directions, and despite the subtraction of power divergences,
one might hope to achieve a better signal-to-noise ratio.
The operator with all the same indices has a power divergence that requires
a complicated nonperturbative procedure to be subtracted and is
not a viable solution. This is unfortunate because the matrix element
of $O_{444}$, for example, does not require any external spatial momentum
and should be optimal in terms of the signal-to-noise ratio.

The mixing pattern becomes more cumbersome for higher dimensional
operators, $n \ge 4$. For $n=4$, choosing
all $4$ indices differently results in the only $2$ irreducible
representations: the totally symmetric and antisymmetric representations.
One then needs a nonvanishing external momentum in all $3$ spatial directions
to calculate the hadronic matrix elements~\cite{Beccarini:1995iv}.
For $n > 4$, mixings are unavoidable.
The complicated renormalization pattern just described,
together with the need of an external momentum to calculate
the hadronic matrix elements, has, {\it de facto}, prevented the lattice calculation
of the moments of parton distribution functions for $n>4$ and rendered the calculations
for $n=3,4$ challenging.

A very interesting idea to circumvent the renormalization issue,
based on the recovery of the continuum rotational symmetry,
was proposed in Ref.~\cite{Davoudi:2012ya}.
Even though the method was not further pursued numerically,
it represents a first important step toward resolving
the theoretical challenges encountered when calculating moments of PDFs.
